% Folder 156-8, batch 01 (archivists' pages 1-20).
% Pass: Opus 5.5 (claude-opus-5-5), 2026-10-02 — first pass, unchecked against the pages by a human.
%
% Pages 4 and 6 are unrelated typed versos (typescript pages of another text,
% scanned upside down) and are not transcribed. Page 1 is the folder's cover,
% inscribed in his hand; its words become the opening heading.
\documentclass[11pt,a4paper]{article}
\input{../preamble/grothendieck}

\folder{156-8}
\batch{1}
\pages{1}{20}
\foldertitle{[Chapitre] VIII. Analysis situs (quatrième mouture) : notes manuscrites (26/06-04/07/1986).}
\dating{1986}
\watermark{Édition de démonstration}

\begin{document}

\section{GF VIII}

\note{intitulé porté de sa main sur la couverture (p.~1), avec, plus bas : « Voir page \uncertain{46} Déf de magasins ». « GF » est le sigle qu'il porte en tête des feuillets de cette série.}

\page{2}
\note{page-tableau : une hiérarchie des types de magasins, reliés par des traits. Le schéma ci-dessous reprend les traits du feuillet (sans flèches : la page ne trace que des traits) ; le contenu de chaque case est transcrit ensuite. En haut à droite : « GF VIII ».}

\begin{tikzcd}[column sep=small, row sep=small, nodes={font=\scriptsize}]
 & \text{Magasins généraux} \arrow[dl, no head] \arrow[dr, no head] & \\
\text{Magasins locaux} \arrow[dr, no head] & & \text{Magasins fidèles} \arrow[dl, no head] \arrow[d, no head] \\
 & \text{Magasins ponctuaires} \arrow[d, no head] & \text{Maquettes} \\
 & \text{Magasins ensemblistes} \arrow[d, no head] & \\
 & \text{Magasins ensemblistes modérés} &
\end{tikzcd}

\textbf{Magasins généraux} : $\mathcal{M}$, $\leq$, $\ll$, $|\circ|$ ; Mag~1 -- Mag~4 (p.~22). $\Sigma_{\mathcal{M}} \subset \mathfrak{P}(\mathcal{M})$.

\textbf{Magasins locaux} : $\mathcal{M}$, $\leq$, $\mathring{\ll}$, $|\circ|_{L}$ ; Magloc~1 -- Magloc~3 (p.~23). ($A, B \in \Sigma_{\mathcal{M}}$, $A \cap B = \emptyset \not\Rightarrow A |\circ| B$)\note{la flèche paraît barrée ; lecture du barrement incertaine.}
$\Sigma_{\mathcal{M}} \subset \mathfrak{P}(L)$ (cette inclusion ne commutant pas \uncertain{néc.} à $A \cap B$ \uncertain{ni à} $\bigcup$ (?), et aucun \struck{\ill{}} $\mathrm{Sup}(A,B)$ \ill{}).

\marginal{NB Un magasin « local » \uncertain{i.e. qui se réfère à} \struck{\ill{}} \uncertain{des voisinages particuliers dévolus aux lieux}.}

\textbf{Magasins} \struck{\ill{}} \add{(à ombres)} \textbf{fidèles} : $\mathcal{M}$, $\leq$, $\mathring{\ll}$ ; Magfid~1 -- Magfid~2 (p.~31).
[$X \mapsto \mathrm{Omb}^{\circ}(X)$ \struck{fonction} \add{fonction} support \emph{fidèle}. $A, B \in \Sigma_{\mathcal{M}}$, $A \cap B = \emptyset \Leftrightarrow A |\circ| B$.]

\textbf{Magasins} \struck{ponctuels} \add{ponctuaires} : $\mathcal{M}$, $\leq$, $\ll$ ; Mag ponc~1-3 (p.~24). $\Sigma_{\mathcal{M}} \simeq \mathfrak{P}(L)$.

\marginal{NB Un lieu $x \in L$ est dit « ponctuel » \add{ou « un point » de $\mathcal{M}$}, si $\forall X \in \mathcal{M}$, $x |\circ| X \Leftrightarrow x \notin \mathrm{Omb}^{\circ}(X)$ (cela \uncertain{entraîne}, si $y \in \mathcal{M}$, $x |\circ| y \Leftrightarrow x \neq y$). Les mag. « ponctuels » \uncertain{sont} locaux dont tous les lieux sont \struck{\ill{}} ponctuels. Ce sont les seuls magasins qui \struck{\ill{}} correspondent \uncertain{plus ou moins} communément \uncertain{admis} pour « la topologie ».}

\textbf{Maquettes} : $\mathcal{M}$, $\leq$ ; Maq~1 : $\leq$ rel. d'ordre (p.~25). $\Sigma_{\mathcal{M}} \simeq \mathfrak{P}(\mathcal{M})$.
(NB Une maquette \add{\uncertain{ultra}} \uncertain{associée à} un espace \struck{top.} satisfaisant l'axiome de séparation $T_0$, n'est un magasin \struck{\ill{}} ponctuaire que si elle est discrète.)

\textbf{Magasins ensemblistes} : $L$, $\mathcal{M} \subset \mathrm{Figél}(L)$ ; Magens~1, Magens~2 (p.~9).

\textbf{Magasins} \add{ensemblistes} \textbf{modérés} : $L$, $\mathcal{M} \subset \mathrm{Figél}(L)$ ; Magens~1 -- Magens~3.

\marginal{Correspondent à des « topologies modérées », qui n'ont été développées à présent que pour la seule structure linéaire par morceaux \struck{\ill{}} \uncertain{sur} $\mathbb{R}$.}

\marginal{NB Il y a des magasins modérés non ensemblistes, p.ex. les maquettes.}

NB Je préfère finalement, dans $\mathcal{M}$, noter $\trianglelefteq$ la relation d'incidence, et noter \struck{\ill{}} \uncertain{de même} la relation d'incidence entre $\mathcal{M}$ et $\mathcal{F}$. La notation $\lhd$ sera l'incidence stricte. La relation d'inclusion entre figures sera notée \struck{\ill{}} $\subseteq$, \uncertain{où il y a} confusion à craindre, \struck{\ill{}} (inclusion des figures). Je noterai $\overset{s}{\trianglelefteq}$ pour les subdivisions ?\note{les symboles de cette dernière note sont surchargés ; lecture de $\overset{s}{\trianglelefteq}$ incertaine.}

\page{3}
\section{Multistructures orientées -- magasins orientés}

Données d'un magasin \struck{\ill{}} orienté : \add{a :}
\[
  \mathcal{M},\ \underset{\text{incidence immédiate}}{\lhd},\ \mathring{\ll},\ |\circ|,\
  \underset{\text{involution sans pt fixe}}{X \mapsto X^{-}},\ \mathcal{M}_0^{+}
\]
\note{sous $X \mapsto X^{-}$ est ajouté : « multistr. orientée opposée » ; après $\mathcal{M}_0^{+}$ un mot biffé.}

\emph{Magor.~1} a) $\lhd$ antiréflexive. Le préordre strict engendré $\leq$ ($X < Y$ ssi $\exists\, X_0 = X \lhd X_1 \lhd \cdots \lhd X_{n-1} \lhd X_n = Y$) est un ordre strict (pas de chaînes fermées, i.e.\ on ne peut avoir $X < X$). Il satisfait la condition des chaînes, \struck{\ill{} tout élément est de dim finie, et les éléments minimaux \ill{}} (en particulier, tout $X \in \mathcal{M}$ de dim finie pour $\leq$).

\marginal{\uncertain{(b) donnée de $\lhd$} ainsi revient à donner une relation d'ordre satisfaisant la cond. des chaînes descendantes.}

b) $\mathring{\ll}$ relation d'ordre.

c) $|\circ|$ \struck{symétrique} antiréflexive.

d) $X \mapsto X^{-}$, involution sans pt fixe.

e) $\mathcal{M}_0^{+}$ est une partie de $\mathcal{M}$.

(\emph{Magor 1'} \struck{a)}, (a) $X \lhd Y \Rightarrow X^{-} \lhd Y^{-}$,

b) $X \mathring{\ll} Y \Rightarrow X^{-} \mathring{\ll} Y^{-}$

c) $X |\circ| Y \Rightarrow X^{-} |\circ| Y$.

\marginal{NB On n'a pas \ill{} « \ill{} $X \lhd Y$ » \uncertain{ou} $X \lhd Y'$.}
\drawing{en marge de Magor~1', une petite boucle d'un point $x$ à $Y$.}

\emph{Magor~2} a) Si $X, Y$ tels que $\exists\, Y'$, $X \mathring{\ll} Y' \leq Y$, alors $Y'$ \add{unique}.

b) Si $X \mathring{\ll} Y$, et $X' \lhd X$, $\exists\, Y' \lhd Y$ \ill{} \add{$X' \ll Y'$}.

\emph{Magor~3} Si $X |\circ| Y$, \struck{et} $X' \mathring{\ll} X$, $Y' \mathring{\ll} Y$, \ill{} $X' |\circ| Y'$.

\emph{Magor~4} Si $X < Y$, alors $X |\circ| Y$.

\emph{Magor~6}\note{une flèche reporte Magor~6 et Magor~5 avant Magor~4 ; son point d'arrivée exact est incertain.} a) Soient \struck{$X$} $Z < X$, $\mathrm{codim}(Z,X) = 2$. Alors il n'existe \uncertain{qu'}\emph{un seul} $Y^{\sharp}$, avec $Z \lhd Y^{\sharp} \lhd X$, \uncertain{et de plus il existe} $Y^{\flat}$ avec $Z^{-} \lhd Y^{\flat} \lhd X^{-}$.

b) Soit $X$ avec $\dim X = 1$. Alors il \struck{existe} \uncertain{existe} un seul $x \in \mathcal{M}_0^{+}$, $x \lhd X$ (\uncertain{i.e.\ $x = \mathcal{M}_0$}).

\emph{Magor~5} a) $x \in \mathcal{M}_0^{+} \Rightarrow \mathrm{codim}\, x = 0$, i.e.\ $\nexists\, y$ avec $y \lhd x$.

b) On a $\mathcal{M}_0 = \mathcal{M}_0^{+} \amalg a(\mathcal{M}_0^{+})$, i.e.\ $x \in \mathcal{M}_0^{+} \Rightarrow x^{-} \notin \mathcal{M}_0^{+}$, et $\forall x \in \mathcal{M}_0$, on a $x \in \mathcal{M}_0^{+}$ ou $x^{-} \in \mathcal{M}_0^{+}$.

\page{5}
\emph{Maquettes orientées} : quand $X \mathring{\ll} Y \Rightarrow X = Y$, $X |\circ| Y \Leftrightarrow X \neq Y$.

\emph{Ex} Les \uncertain{systèmes} de $\mathcal{M}$ sont de dim $\leq 1$, $\mathring{\ll} = \mathrm{id}$, $x |\circ| y \Leftrightarrow x \neq y$.
\struck{$\mathcal{M} = \mathcal{M}_0^{+}$}
On pose $\mathcal{M}_0^{+} = S$ (« sommets »), $\mathcal{M}_1 = A$ (« arêtes orientées »).
On a involution sans points fixes $a$ sur $A$.
\struck{\ill{}} Magor~5 \ill{} $A \xrightarrow{\;o\;} S$ [\ill{} application].

Les maquettes orientées \struck{\ill{}} de dim $\leq 1$ \uncertain{$\lhd$} (« les graphes »).

\emph{NB} Cette structure \ill{} \uncertain{par} une maquette, \ill{} \uncertain{arêtes} \struck{\ill{}} \ill{} des maquettes.

\emph{Cas de \struck{\ill{}} dim 2} \ill{}, on trouve la description des \ill{} cellulaires combinatoires de dim $\leq 2$.

\note{la suite de la mouture orientée ne vient pas ; la page 7 reprend tout à neuf.}

\page{7}
\note{la numérotation des archivistes portée au crayon saute le 7 : ce feuillet porte « 8 », et le décalage d'une unité se maintient jusqu'au dernier feuillet du lot, qui porte « 21 ». On suit ici le rang du feuillet dans le lot. Il numérote « 1 » ; en haut à droite, « GF VIII ». En marge : « 26 juin 86 ».}
\section{Analysis Situs (quatrième mouture)}

\textbf{I} \emph{Récapitulation en termes de « magasins ».}
Je reprends le formalisme algébrique d'une « algèbre en multistructures », ou « magasin » (de multistructures, magasin « fournisseur » \uncertain{de} « \ill{} »). La structure est celle d'un ens.\ $\mathcal{M}$ avec trois relations
\[
  \leq,\quad \ll,\quad |\circ|
\]
satisfaisant (cf GF~VII, p.~110) :

\emph{Mag~1} a) $X \leq Y$ et $X \ll Y$ sont des relations d'ordre, et b)
\[
  X \leq Y \Longrightarrow X \ll Y .
\]
Si $X \in \mathcal{M}$, on pose $\widetilde{X} = \mathcal{M}_{\leq X} = \{ Y \in \mathcal{M} \mid Y \leq X \}$.

\emph{Mag~2} $\forall X, Y$ avec $X \ll Y$, l'ens
\[
  \widetilde{Y}^{X} = \{ Z \in \widetilde{Y} \mid X \ll Z \}
\]
admet un plus petit élément. On écrit $X \mathring{\ll} Y$ quand \uncertain{cet élément} \add{minimal} est $Y$ lui-même, i.e.\ $\widetilde{Y}^{X} = \{Y\}$.

\emph{Mag~3} ⓐ $X |\circ| Y$ est symétrique et antiréflexive, i.e.\ $X |\circ| Y$ implique $Y |\circ| X$ et $X \neq Y$. ⓑ De plus
\[
  X |\circ| Y,\ X' \mathring{\ll} X,\ Y' \mathring{\ll} Y \Longrightarrow X' |\circ| Y' .
\]

\emph{Mag~4} Si $Y, Z \leq X$, \struck{alors} $Y \neq Z$, on a $Y |\circ| Z$.

\marginal{NB Cela \uncertain{assouplit légèrement} l'ancien \uncertain{At~2}.}

\emph{Définitions} Pour $X, Y \in \mathcal{M}$, on pose
\note{il écrit une relation formée de $\leq$ et $\geq$ superposés ; on la note ici $\lessgtr$.}
\[
\begin{aligned}
  X \lessgtr Y &\overset{\text{déf}}{\Longleftrightarrow} \forall X' \in \widetilde{X} \smallsetminus \widetilde{X} \cap \widetilde{Y},\ Y' \in \widetilde{Y} \smallsetminus \widetilde{X} \cap \widetilde{Y},\ \text{on a } X' |\circ| Y' \\
  X \parallel Y &\overset{\text{déf}}{\Longleftrightarrow} \forall X' \in \widetilde{X},\ Y' \in \widetilde{Y},\ \text{on a } X' |\circ| Y' .
\end{aligned}
\]
Donc
\[
  X \parallel Y \Longleftrightarrow X \lessgtr Y,\ \text{et } \widetilde{X} \cap \widetilde{Y} = \emptyset
\]
(i.e.\ $\{X, Y\}$ non minoré pour $\leq$).

\page{8}
\note{il numérote « 2 », avec en tête « (inconditionnelle) », lecture incertaine.}
On appelle \emph{figures} (du magasin, un élément de l'ens)
\[
  \mathcal{F}_{\mathcal{M}} = \{ \mathfrak{F} \subset \mathcal{M} \mid \mathfrak{F} \text{ fermé dans } \mathcal{M} \text{ pour } \leq, \text{ et } \forall X, Y \in \mathfrak{F},\ \text{on a } X \lessgtr Y \}
\]
L'ens $\mathcal{F}_{\mathcal{M}}$ est l'ensemble des « figures \uncertain{inconditionnelles} » (ou « figures » tout court) du magasin. On sera amené, à l'occasion, de faire sur les figures envisagées deux types possibles de restriction :

1°) Conditions du type « finitude » ou « local finitude ».

\struck{\ill{}} 2°) Condition \struck{restrictive} « polyédrale » : si $X, Y \in \mathfrak{F}$, et $\widetilde{X} \cap \widetilde{Y} \neq \emptyset$, on \struck{veut} veut qu'il \struck{$\widetilde{X} \cap \widetilde{Y}$} admette un plus grand élément, i.e.\ soit de la forme $\widetilde{Z}$, ou encore que $\mathrm{Inf}(X,Y)$ existe.

\marginal{Voir pages 1bis, 2bis, à mettre \add{avant} les sorites sur figures.}

\marginal{$\mathcal{F}_{\mathcal{M}}$ sera appelé \emph{\uncertain{l'édition maximale}} associé\supplied{e} au magasin $\mathcal{M}$.}

On étend : les relations d'ordre $\leq$ et $\ll$ de $\mathcal{M}$ (en plongeant $\mathcal{M}$ dans $\mathcal{F}$ via les « figures élémentaires » $\widetilde{X}$, $X \in \mathcal{M}$
\[
  \mathcal{M} \hookrightarrow \mathcal{F}_{\mathcal{M}},\qquad X \mapsto \widetilde{X} = \mathcal{M}_{\leq X}
\]
\uncertain{est injective} et permet \uncertain{de l'identifier} \ill{}, \ill{} $\mathcal{F}_{\mathcal{M}}$ \ill{} figures élémentaires, $X \in \mathcal{M}$ : la figure élémentaire \uncertain{associée} \ill{})
\[
  F \leq G \overset{\text{déf}}{\Longleftrightarrow} F \subset G
\]
NB Le Sup dans $\mathcal{F}$ pour $\leq$ est $\bigcup$ dans $\mathfrak{P}(\mathcal{M})$, le Inf $\bigcap$.
\[
\begin{aligned}
  F \ll G &\overset{\text{déf}}{\Longleftrightarrow} \forall X \in \widetilde{F},\ \exists\, Y \in \widetilde{G},\ \text{avec } X \ll Y \\
  F \lessgtr G &\Longleftrightarrow \forall X \in \widetilde{F},\ Y \in \widetilde{G},\ \text{on a } X \lessgtr Y \\
  &\Longleftrightarrow \{F, G\} \text{ majoré dans } \mathcal{F} \\
  &\Longleftrightarrow \mathrm{Sup}(F,G) \text{ existe dans } \mathcal{F} \\
  &\Longleftrightarrow F \cup G \in \mathcal{F}
\end{aligned}
\]
\note{après « existe dans $\mathcal{F}$ », un mot biffé.}

\marginal{\ill{} on identifie les figures \uncertain{d'une figure} de $\mathcal{F}_{\mathcal{M}}$, et pour une figure $F$ de $\mathcal{F}_{\mathcal{M}}$ \ill{} \uncertain{à une partie}, on identifie $F$ à $\widetilde{F} = \{ X \in \mathcal{M} (\subset \mathcal{F}) \mid X \leq F \}$.}

\page{9}
\note{il numérote « 1bis » ; en marge : « 27 juin ».}
Avant de passer aux figures, il faudrait donner quelques compléments sur $\mathcal{M}$ lui-même.

\struck{\emph{Prop} Si $Y \mathring{\ll} X$ et $Y \mathring{\ll} X'$, alors $X \lessgtr X'$.}

\struck{\emph{Prop} Supposons $X \lessgtr Y$. Si $Z \mathring{\ll} X$, alors}
\[
  \text{\struck{$Z \ll Y \Longleftrightarrow X \ll Y$.}}
\]
\struck{\emph{Dém} $\Leftarrow$ trivial, prouvons $\Rightarrow$. Par Mag~2, $\exists\, Y'$ avec $Z \mathring{\ll} Y' \leq Y$}

\struck{\emph{Prop~1} Si $X' \leq X$, $Y' \leq Y$, alors $X \lessgtr Y \Rightarrow X' \lessgtr Y'$}
\note{ces trois débuts sont biffés de traits obliques ; le dernier est encadré.}

\emph{Prop~$1_0$} a) $X \lessgtr Y$ est symétrique et réflexive

b) Si $X' \leq X$, $Y' \leq Y$, alors $X \lessgtr Y \Rightarrow X' \lessgtr Y'$.

\emph{Dém.} a) \struck{symétrique} triviale par \emph{Mag~3} ⓐ.

b) Soit $X'' \in \widetilde{X'} \smallsetminus \widetilde{X'} \cap \widetilde{Y'}$, $Y'' \in \widetilde{Y'} \smallsetminus \widetilde{X'} \cap \widetilde{Y'}$. \add{Prouvons $X'' |\circ| Y''$.} Si $X'', Y'' \notin \widetilde{X} \cap \widetilde{Y}$, alors $X'' |\circ| Y''$ par déf de $X \lessgtr Y$. Si $X'' \in \widetilde{X} \cap \widetilde{Y}$, on a $X'' \in \widetilde{Y}$, $Y'' \in \widetilde{Y'} \subset \widetilde{Y}$, d'autre part évidemment $X'' \neq Y''$, donc $X'' |\circ| Y''$ par \emph{Mag~4}.

\emph{\struck{Prop.} \add{Cor~1}} Soient \struck{$X$}$X, Y$ avec $X \lessgtr Y$. Soit $Z$ avec $Z \mathring{\ll} X$. Alors
\[
  Z \ll Y \Longleftrightarrow \text{\struck{$Z$}}\, X \ll Y
\]
Énoncé \ill{} \uncertain{se déduit} corollaire

\emph{Corollaire~2} Si $Z \mathring{\ll} X$, $Z \mathring{\ll} Y$, et si $X \lessgtr Y$, alors $X = Y$.

En effet, \ill{} $X \ll Y$ et $Y \ll X$, donc $X = Y$ par Mag~1 a).

Prouvons \struck{Prop.~2} \add{Cor~1} $\Leftarrow$ évident, prouvons
\[
  Z \mathring{\ll} X,\ Z \ll Y \Longrightarrow X \ll Y \qquad (\text{si } X \lessgtr Y).
\]
\struck{Soit} Par \emph{Mag~2}, $\exists\, Y'$ avec $Z \mathring{\ll} Y' \leq Y$, par prop~$1_0$ on a $X \lessgtr Y'$, il suffit de prouver $X = Y'$ (d'où $X \leq Y$), i.e.\ \ill{} \uncertain{Mais si on avait} $X \neq Y'$, on se ramène au cas \ill{}
\[
  \text{\struck{$Z \mathring{\ll} X,\ Z \mathring{\ll} Y,\ X \lessgtr Y \Longrightarrow X = Y$.}}
\]
\struck{Prouvons $Z \in \widetilde{X} \cap \widetilde{Y}$, on \ill{} $X \leq Y$ (car \ill{} $Y \leq X$} \struck{Il suffit de prouver $\widetilde{X} \subset \widetilde{Y}$, sinon on aurait $X \parallel Y$ (cf lemme), \ill{} $Z |\circ| Z$} \struck{Sinon, on aurait $Z \mathring{\ll} X$, $Z \mathring{\ll} Y$ \ill{} Mag~3 b), absurde par Mag~3 a)}, on aurait par prop~2 $X |\circ| Y$, donc $Z |\circ| Z$ par Mag~3 b), absurde par Mag~3 a).
\note{le bas de la page est un palimpseste de tentatives biffées de traits obliques ; seule la dernière ligne n'est pas biffée.}

\marginal{\emph{Prop.~2} Si \uncertain{$X \neq Y$}, $X |\circ| Y \Leftrightarrow X \lessgtr Y$. \emph{Dém} $\Rightarrow$ par Mag~3 a). Inv. Si \uncertain{$X \neq Y$}, $X \lessgtr Y$, \struck{et} si \uncertain{$X \not\leq Y$ et $Y \not\leq X$}, i.e.\ $X \notin \widetilde{X} \cap \widetilde{Y}$, $Y \notin \widetilde{X} \cap \widetilde{Y}$, \ill{} $X |\circ| Y$ par déf.\ de $X \lessgtr Y$. Si $X \leq Y$, \ill{}}

\page{10}
\note{il numérote « 2bis ».}
\emph{Cor.~3} Soit $X \lessgtr Y$. Alors
\[
  X \mathring{\ll} Y \Longleftrightarrow X \leq Y
\]
\struck{En effet,} On applique cor~1 avec $Z = X$.

\emph{Proposition~3.} \struck{Supposons} $X \mathring{\ll} Y \mathring{\ll} Z$ \struck{Alors} \add{implique} $X \mathring{\ll} Z$. \struck{$\Longleftrightarrow X \mathring{\ll} Y$ et $Y \mathring{\ll} Z$}

\struck{Supposons $X \mathring{\ll} Y$ et $Y \mathring{\ll} Z$, prouvons $X \mathring{\ll} Z$, i.e.\ $X \ll Z' \leq Z \Rightarrow Z' = Z$. Comme}

Par Mag~2 $\exists\, Z'$ avec $X \mathring{\ll} Z' \leq Z$, \struck{il faut} \struck{par Mag} prouvons $Z' = Z$. Sinon, par Mag~4 on aurait $Z' |\circ| Z$, et par Mag~3 b) $X |\circ| Z$, d'autre part par Mag~3 b) cela implique $X |\circ| Y$ puisque $Y \mathring{\ll} Z$, puis $X |\circ| X$ puisque $X \mathring{\ll} Y$, \struck{d'où $X |$} mais $X |\circ| X$ contredit Mag~3 a).

\note{à gauche, un petit schéma : $X \mathring{\ll} Y \mathring{\ll} Z$, avec $Z'$ au-dessous, relié à $X$ par $\mathring{\ll}$ et à $Z$ par $\leq$.}

\emph{NB} \struck{Si} Supposons \add{$X \ll Y \ll Z$ et} $X \mathring{\ll} Z$. Alors $Y \mathring{\ll} Z$, car si on avait $Y \ll Z' < Z$, \uncertain{on aurait} $X \ll Z' < Z$, contrairement à $X \mathring{\ll} Z$. \struck{Prouvons} Mais on n'a pas \uncertain{nécess.} $X \mathring{\ll} Y$, exemple :
\drawing{un segment gradué ; une accolade inférieure désigne le segment entier, $Z$ ; une accolade supérieure en désigne une partie, $Y$ ; un point marqué à l'extrémité de cette partie, désigné par une flèche, est $X$.}

\emph{Cas des figures}

\emph{Prop~$1_F$} a) $F \lessgtr G$ sym.\ réfl.

b) $F \lessgtr G$, $F' \leq F$, $G' \leq G \Rightarrow F' \lessgtr G'$

\emph{Prop~$2_F$} Si $F \lessgtr G$, \struck{Prop} alors $F \ll G \Leftrightarrow F \leq G$

\emph{Cor} Dans l'ens \uncertain{$\mathcal{S} \subset \mathcal{F}_{0}(F)$}, $\ll$ et $\leq$ coïncident.

\page{11}
\note{il numérote « 3 » : la suite de la page 8 (« 2 »).}
\[
\begin{aligned}
  F \parallel G &\Longleftrightarrow \forall X \in \widetilde{F},\ Y \in \widetilde{G},\ \text{on a } X \parallel Y \\
  &\Longleftrightarrow F \lessgtr G \text{ et } F \cap G = \emptyset
\end{aligned}
\]
On n'écrit $F |\circ| G$ que si $F, G$ sont élémentaires, i.e.\ sont « $\in \mathcal{M}$ », idem pour $\mathring{\ll}$.

\struck{\emph{Proposition} Soient $F, G \in \mathcal{F}$, $L = F \cap G$, $F' \ll F$, $G' \ll G$,} \struck{On a}

\emph{Proposition~1} Propriétés \uncertain{usuelles} de $\mathcal{F}, \leq$ : a) \uncertain{stabilité} \ill{} et b) \uncertain{existence} \add{« division » des parties $\mathfrak{F}$} \ill{} ensemble de parties de $\mathcal{M}$.

\marginal{y compris l'existence \uncertain{des Sup} \ill{} \uncertain{des Inf} ($\mathfrak{F} \neq \emptyset$)}

\emph{Proposition~2} a) $F \lessgtr G$, $F' \leq F$, $G' \leq G \Rightarrow F' \lessgtr G'$

\struck{b) $F \parallel G$, $F' \ll F$, $G' \ll G \Longrightarrow F' \parallel G'$}

b) \struck{En} Si $F = \bigcup_i F_i$ \uncertain{Sup} \ill{} figures,
\[
  \text{alors}\quad G \lessgtr F \;\text{\struck{$\Longleftrightarrow$}}\; \forall i,\ G \lessgtr F_i .
\]
Donc si $F = \bigcup F_i$, $G = \bigcup G_j$, $F \lessgtr G$ ssi $\forall i, j$, $F_i \lessgtr G_j$.

\marginal{Prop~1, Prop~2 \uncertain{concernent seulement} $\leq$}

\emph{Proposition~3} \struck{Soient} \add{Soient} $F, G$, $F \lessgtr G$, \struck{$F' \ll F$,} \struck{$L = F \cap G$, $F' \ll F$, $G' \ll G$,}

\note{il écrit les inégalités verticalement ; le signe $\geq$ placé entre deux lignes se lit de haut en bas.}
\[
  \begin{array}{ccc}
    G & \ll & F \\
      &     & \geq \\
      &     & F'
  \end{array}
\]

alors $\exists\, G'$ avec \struck{$\widetilde{G} \cap \mathrm{Omb}(F')$}

\[
  \begin{array}{ccc}
    G  & \ll & F \\
    \geq &   & \geq \\
    G' & \ll & F'
  \end{array}
\]

$G' = \mathrm{Inf}^{\ll}(F', G)$.

On a
\[
  \widetilde{G'} = \widetilde{G} \cap \mathrm{Omb}(F')
\]
Posant $G' = G_{F'}$ \add{i.e.\ $G(F')$}, ou $F'_{G}$ \add{i.e.\ $F'|G$}, on a
\[
  F' \longmapsto G_{F'} : \mathrm{Ssfig}(F) \longrightarrow \mathrm{Ssfig}(G)
\]
\struck{est croissante et} \add{commutant aux} Sup et aux Inf quelconques.

\page{12}
\note{il numérote « 4 ».}
\emph{Cor.~1} Si $G \leq F$, alors $G' = G \wedge F'$ \add{(donc $\widetilde{G'} = \widetilde{G} \cap \widetilde{F'}$), i.e.}
\[
  \mathrm{Inf}^{\ll}(G, F') = \mathrm{Inf}^{\leq}(G, F')
\]
\emph{Cor.~2} (\uncertain{cf lemme p.~5}) $L = F \cap G$,

\emph{Prop~5} Soient $F, G$ avec $F \lessgtr G$, $F' \ll F$, $G' \ll G$, d'où $F'_{L} \ll L$, $G'_{L} \ll L$. Alors
\struck{\ill{}}
\[
\begin{aligned}
  \text{a)}\quad & F' \lessgtr G' \Longleftrightarrow F'_{L} \lessgtr G'_{L} \\
  \text{b)}\quad & F' \parallel G' \Longleftrightarrow F'_{L} \parallel G'_{L}
\end{aligned}
\]
\emph{Cor.} Si $F'_{L} = \emptyset$, alors $F' \parallel G'$ pour tout $G' \leq G$.

\struck{(en fait $F'' \parallel G$)}

Prenant $L = \emptyset$, on trouve

\emph{Cor} \struck{$F \parallel G$,} Si $F' \ll F$, $G' \ll G$, alors
\[
  F \parallel G \Longrightarrow F' \parallel G'
\]
\note{les énoncés depuis Prop~5 sont réunis par une accolade ; en marge, une flèche descendante « \uncertain{demande prop.~4} ».}

\emph{Prop.~4} (Recollement des raffinements)

Soit $F = \bigcup_i F_i$. Alors
\[
\begin{array}{ccccc}
  \mathrm{Raff}(F) & \longrightarrow & \prod_i \mathrm{Raff}(F_i) & \rightrightarrows & \prod_{i,j} \mathrm{Raff}(F_i \cap F_j) \\
  G & \longmapsto & (G|F_i)_i & &
\end{array}
\]
exact.

\marginal{On pose $\mathrm{Raff}(F) = \{ F' \ll F \}$ \ill{} $G|F \ll F$ \ill{}}

\emph{Cor.~1} Soit $F \in \mathcal{F}$. Alors
\[
  \mathrm{Raff}(F) \longrightarrow \prod_{X \in \widetilde{F}} \mathrm{Raff}(X)
\]
\struck{\ill{}} est injectif, et l'image est formée des familles $(G_X)_{X \in \widetilde{F}}$ de raffinements telles que $Y \leq X$ implique $G_Y = G_X | Y$. On a donc
\[
  \mathrm{Raff}\, F = \varprojlim_{X \in \widetilde{F}} \mathrm{Raff}(X)
\]

\page{13}
\note{il numérote « 5 ».}
\struck{Tout ceci est récapitulatif, \uncertain{on fait pas} \ill{} \ill{} \uncertain{directement} $|\circ|$.}
\note{ces deux lignes sont barrées de grandes croix ; elles sont reprises plus bas.}

\struck{\emph{Prop.~4}} \emph{Cor~2} Soit $G \ll F = \bigcup F_i$. Alors $G \leq F$ ssi \struck{Prop. pour} $\forall i$, $G|F_i \leq F_i$. \struck{\ill{}} Si $G, G' \ll F$, alors $G \lessgtr G'$ ssi $\forall i$, $G_{F_i} \lessgtr G'_{F_i}$, idem pour $\parallel$, $\leq$, $\ll$ …

\marginal{devrait \uncertain{venir} \uncertain{après} prop.~3.}

Tout ceci est récapitulatif, et ne fait pas intervenir explicitement $|\circ|$, si ce n'est par l'intermédiaire de $\lessgtr$. Il faudrait, dans ce contexte, reprendre aussi rapidement \add{\uncertain{les}}
\begin{enumerate}
  \item[1)] Les axiomes relatifs aux lieux, ombres, multiombres.
  \item[2)] Le \uncertain{lien} avec le cas ensembliste et quasi-ensembliste (où les lieux jouent justement un rôle prépondérant).
\end{enumerate}
Mais surtout
\begin{enumerate}
  \item[3)] Revoir la théorie des \uncertain{subdivisions}.
\end{enumerate}
\struck{Comme} Également
\begin{enumerate}
  \item[4)] Voir ce que devient l'axiomatique $\leq$, $\ll$, $|\circ|$, dans le cas (le plus courant sûrement) où $\ll$ peut s'exprimer en termes de $\leq$, $|\circ|$. Mais pour les formules, nous aurons besoin de donner un sens, pour $X \in \mathcal{M}$, à
  « $X^{\circ} = X \smallsetminus \partial X$ »,
  et pour $X, Y \in \mathcal{M}$, à « l'inclusion »
  « $X^{\circ} \subset Y^{\circ}$ »
  i.e.\ il faut d'abord passer par 3.
\end{enumerate}
Traitons quand même 1° et 2°).

\struck{Les axiomes pertinents devraient}

\page{14}
\note{il numérote « 6 ».}
Axiomes pour les lieux, i.e.\ les éléments de $\mathcal{M}$ minimaux pour $\ll$. Pour $x, y \in L$,
\[
  \underset{\text{ens.\ des lieux}}{L \subset \mathcal{M}},
\]
il est tautologique que
\[
  x \parallel y \Longleftrightarrow x |\circ| y .
\]
\struck{On a aussi} \struck{Notons aussi,} On pose comme d'habitude, si $X \in \mathcal{M}$
\[
\begin{aligned}
  \mathrm{omb}(X) &= \{ x \in L \mid x \ll X \} = \mathrm{Omb}(X) \cap L \\
  \mathrm{omb}(F) &= \textstyle\bigcup_{X \in F} \mathrm{omb}(X) \qquad \text{pour } F \subset \mathcal{M} \text{ (p.ex.\ } F \in \mathcal{F}) \\
  \text{\struck{$\mathrm{omb}(X)^{\circ}$}} &\;\text{\struck{$= \mathrm{omb}(X)$}} \\
  \mathrm{omb}(X)^{\circ} &= \mathrm{omb}(X) - \mathrm{omb}(\partial X)
\end{aligned}
\]
\marginal{plus gén.\ (en regard de la ligne $\mathrm{omb}(F)$)}
où
\[
  \partial X = \widetilde{X} \smallsetminus \{X\} \;\text{\struck{$F$}}\; = \{ Y \in \mathcal{M} \mid Y \lneq X \} \in \mathcal{F}
\]
\note{le $\lneq$ final est lu sur un $\leq$ surchargé ; lecture incertaine.}
Notons
\[
  x \in \mathrm{omb}(X)^{\circ} \Longleftrightarrow x \mathring{\ll} X
\]
On a alors, par Mag~3
\[
  X \parallel Y \Longrightarrow \forall x \in \mathrm{omb}(X)^{\circ},\ y \in \mathrm{omb}(Y)^{\circ},\ \text{on a } x \parallel y .
\]
Ceci n'est pas bien utile, quand $\mathrm{omb}(X)^{\circ}$ ou $\mathrm{omb}(Y)^{\circ}$ est $\emptyset$.

\emph{Mag$L$~1} $\forall X \in \mathcal{M}$, $\mathrm{omb}(X)^{\circ} \neq \emptyset$

Ceci nous \add{permet de définir, et} donne des propriétés agréables pour
\[
\begin{aligned}
  \mathcal{M} &\longrightarrow \mathrm{Figél}(L) &\quad&\text{ou}\quad& \mathcal{F} &\longrightarrow \mathrm{Fig}(L) \\
  X &\longmapsto \mathrm{mulomb}(X) &&\text{ou}& F &\longmapsto \mathrm{mulomb}(F)
\end{aligned}
\]
savoir : croissant \struck{compat.} pour $\leq$, $\ll$, commutant \add{aux \struck{\ill{}} Sup \struck{quelconques} \uncertain{finis}}, compatible avec \add{$\leq$,} $|\circ|$, $\parallel$ … Il faudrait revoir la notion d'homomorphisme et de \uncertain{plongement} pour les \struck{ateliers} magasins (\uncertain{ou} pour les ateliers, par \uncertain{exemple}) cf GF~VII p.~54 et GF~VI …

\page{15}
\note{il numérote « 7 ».}
\emph{Mag$L$~2\add{\uncertain{bis}}} Soient $X, Y \in \mathcal{M}$, \struck{tels donc} tels que
\[
  \forall x \in \mathrm{omb}(X)^{\circ},\ y \in \mathrm{omb}(Y)^{\circ},\ \text{on a } x \parallel y .
\]
Alors $X |\circ| Y$.

\marginal{\uncertain{implique} Mag$L$~1 (à cause de l'antiréflexivité de $|\circ|$)}

Cela signifie donc que $|\circ|$ sur $\mathcal{M}$ se déduit de la \struck{\ill{}} relation $\parallel$ sur $L$ \add{dans les cas pratiques}. Il semble que chaque fois qu'on a Mag$L$~1, on a aussi Mag$L$~2, qui semble donc l'axiome le plus utile.

\emph{Proposition -- Scholie.} Se donner \add{sur $\mathcal{M}$} \add{structure de} une \uncertain{magasin} « strictement divisible » (i.e.\ satisfaisant à Mag$L$~2, \ill{}) revient au même que de se donner $\leq$, $\ll$ satisfaisant Mag~1, Mag~2 p.~1, \add{(la condition Mag$L$~1 (p.~6)} \emph{plus} une relation
\[
  x \parallel y \quad \text{relation \emph{sur $L$}}
\]
satisfaisant les seules conditions

\emph{Mag div~3} La relation $x \parallel y$ sur $L$ est symétrique et antiréflexive.

\emph{Mag div~4} Si $Y, Z \leq X$ dans $\mathcal{M}$, $Y \neq Z$, et si $y \in \mathrm{omb}(Y)^{\circ}$, $z \in \mathrm{omb}(Z)^{\circ}$, alors $y \parallel z$.

En effet, dans ce cas on \emph{définit} $|\circ|$ dans $\mathcal{M}$ par la condition \uncertain{envers} Mag$L$~2, et on vérifie \uncertain{immédiatement} que \struck{cette} celle-ci satisfait Mag~3\add{*}, Mag~4, et \uncertain{réciproquement}

\marginal{\uncertain{mais pour la réciproque} il faut \uncertain{que} $\mathrm{omb}(X)^{\circ} \neq \emptyset$}

\page{16}
\note{il numérote « 8 » ; la phrase de la page précédente se poursuit.}
la relation \uncertain{se déduit} \uncertain{de sa restriction} à $L$.

Le cas « quasi-ensembliste » est celui où on suppose

\emph{Mag quens} Pour $x, y \in L$, $x \neq y \Rightarrow x \parallel y$.

\emph{Proposition -- Scholie} \struck{Les} \add{Structures de} magasins \struck{\ill{}} \add{pré}-ensemblistes sur $\mathcal{M}$ correspondent aux couples $(\leq, \ll)$ de deux relations d'ordre sur $\mathcal{M}$ sans plus, satisfaisant les seules conditions

\emph{Mag~1} $X \leq Y \Rightarrow X \ll Y$

\emph{Mag~2} $\forall X \ll Y$, l'ens.\ $\widetilde{Y}^{X} = \{ Y' \in \widetilde{Y} \mid X \ll Y' \}$ a un plus petit élément

\emph{Mag$L$~1} $\forall X \in \mathcal{M}$, $\mathrm{omb}(X)^{\circ} \neq \emptyset$, i.e.\ $\exists\, x \in L$, avec $x \mathring{\ll} X$ i.e.\ \struck{$x$} $\widetilde{X}^{x} = \{X\}$.

Voir exemple dans GF~VII p.~57 pour montrer que ceci n'implique pas que $\mathcal{M}$ se plonge dans $\mathrm{Figél}(L)$. \struck{Le magasin \ill{} ensembliste \ill{}}

Pour mémoire : si $\mathcal{M}$ satisfait la condition Mag$L$~1 (sans être nécessairement \uncertain{quasi}-\struck{\ill{}} \uncertain{ici}) alors que \uncertain{ens} \ill{}, on a pour $F \in \mathcal{F}$
\[
\begin{aligned}
  \mathrm{Ssfig}(F) &\xrightarrow{\;\sim\;} \mathrm{Ssfig}(\mathrm{mulomb}(F)) \\
  F' &\longmapsto \mathrm{mulomb}(F')
\end{aligned}
\]
(isom.\ d'ens.\ \uncertain{ordonnés}). Par contre, même dans le cas quasi-ensembliste, on \struck{\ill{}} \uncertain{n'a pas} nécessairement :

\page{17}
\note{il numérote « 9 ». La condition a) annoncée au bas de la page précédente n'est pas écrite ; la page s'ouvre sur b).}
\[
\begin{aligned}
  \text{b)}\quad & \mathrm{mulomb}(F) \ll \mathrm{mulomb}\, G \Longrightarrow F \ll G \\
  \text{c)}\quad & \mathrm{mulomb}(F) \lessgtr \mathrm{mulomb}(G) \Longrightarrow F \lessgtr G .
\end{aligned}
\]
(c'est quand ces \struck{conditions} implications sont satisfaites \add{(en plus de la condition a) ci-dessus)} \uncertain{et il suffit de la poser pour} $F = X$, $G = Y$, $X, Y \in \mathcal{M}$ \add{(spéciaux)} qu'on dit que le magasin est « ensembliste »)

\struck{Il est donc donné par un ensemble de \ill{} telles données \ill{} celles-ci}

\emph{Proposition} -- \emph{Scholie} (La donnée d'un magasin ensembliste, \uncertain{équivaut} à la donnée d'un ens.\ $L$ (l'ens.\ des « lieux » ou « points »), et d'une partie
\[
  \mathcal{M} \subset \mathrm{Figél}(L),
\]
satisfaisant les conditions

\emph{Magens~1} $\forall$ \struck{$F$} $F \in \mathcal{M}$, et tout \struck{$X$} $X \in \widetilde{F}$, l'ens.
\[
  \widetilde{F}_{\leq X} = ( Y \in F \mid Y \subset X ) \subset F \text{ est aussi dans } \mathcal{M}
\]
($\mathcal{M}$ stable par passage aux \ill{} \ill{})

\emph{Magens~2} $\forall x \in L$, $\{\{x\}\} \in \mathcal{M}$.

\ill{} \uncertain{une digression} !!!

\marginal{Donc on a une \uncertain{espèce} de structure $\mathcal{M} \in \mathfrak{P}\mathfrak{P}(L)$, \add{axiomes} \emph{Magens~0} : \uncertain{Tout} $X \in \mathcal{M}$ est une figure ensembliste élémentaire, \ill{} \uncertain{\emph{Magens~2}} \ill{} À peine plus compliqué que la définition d'une « topologie » ! !!!}

\struck{\ill{}} \ill{} de l'axiomatique. En fait, à ce stade de l'axiomatique, \ill{} « ateliers » de « \ill{} évidents », le point de vue des figures nous avait plutôt compliqué la vie. Ici il vaut décidément mieux partir du magasin, et y conserver aussi claire que possible -- quitte à faire ensuite une \add{petite} complication \ill{} les ateliers plus

\page{18}
\note{il numérote « 10 » ; la phrase de la page précédente se poursuit.}
généraux que l'atelier canonique associé à un magasin -- et de revoir les prop essentielles prop~1 à prop~5 (plus d'autres \uncertain{à venir peut-être}) sous cet angle, \uncertain{logiquement} suffisant à présent.

\emph{27 juin} (ou plutôt, 28 juin ! (on ne sait plus …)

\emph{Définition} Soit $\mathcal{F}$, $\leq$ un ensemble ordonné. On appelle \emph{antifiltre} de $\mathcal{F}$ une partie $\mathcal{F}'$ de $\mathcal{F}$ satisfaisant les conditions suivantes

a) $\mathcal{F}'$ est fermé pour $\leq$, i.e.\ $F \in \mathcal{F}'$, $G \leq F \Rightarrow G \in \mathcal{F}'$.

b) Si $F, G \in \mathcal{F}'$, et si $F \vee G$ existe dans $\mathcal{F}$, alors $F \vee G \in \mathcal{F}'$.

\struck{Cela est que les antifiltres \ill{}}
\note{suit un petit encadré biffé, illisible, avec $F$ et $\mathcal{F}$.}

\emph{Définition} Soit $\mathcal{M}$ un magasin. Un \emph{$\mathcal{M}$-atelier}, \add{ou} \emph{atelier strict de $\mathcal{M}$}, est par définition un antifiltre $\mathcal{F}$ (dans l'ens $\mathcal{F}_{\mathcal{M}}$ des $\mathcal{M}$-figures), \struck{tel que} \struck{ce soit vide et} contenant les \struck{figures vides} figures élémentaires \struck{Cette dernière}
\[
  \widetilde{X} \quad (X \in \mathcal{M})
\]
et la figure vide). \struck{condition résulte des autres si $\mathcal{M} \neq \emptyset$.}

\emph{NB} Cette dernière condition résulte des autres si $\mathcal{M} \neq \emptyset$. On voit que les conditions suivantes \uncertain{sont} équivalentes
\begin{enumerate}
  \item[(i)] $\mathcal{M} = \emptyset$
  \item[(ii)] $\mathcal{F} = \{\emptyset\}$
  \item[(iii)] $\mathrm{card}\, \mathcal{F} = 1$
\end{enumerate}
On dit \struck{que} \add{dans ce cas} que \add{$\mathcal{M}$ est} \struck{le} un magasin \struck{\ill{}} \emph{vide} (ou : \struck{le} un \emph{atelier} \struck{$\mathcal{F}$} vide (ou : \uncertain{tout} vide), ce qui est \uncertain{un abus} de langage puisque $\mathcal{F}$ \ill{} n'est pas vide ….

\marginal{\emph{Mag at} •) $F \in \mathcal{F}$, $G \leq F \Rightarrow G \in \mathcal{F}$ \add{(sur $F \subset \mathcal{F}$)} ; b) $F, G \in \mathcal{F}$, $F \lessgtr G \Rightarrow F \cup G \in \mathcal{F}$ ; c) $\forall X \in \mathcal{M}$, $\widetilde{X} \in \mathcal{F}$ ; d) $\mathcal{F} \neq \emptyset$ (i.e.\ $\emptyset \in \mathcal{F}$).}
\marginal{$\mathcal{F}$ (\uncertain{contenu dans} $\mathfrak{P}(\mathcal{M})$) \uncertain{\emph{commun}} \ill{} $\mathcal{M}$ -- \uncertain{vérifie que} $\Leftrightarrow$ c)}

\page{19}
\note{il numérote « 11 ».}
Les \struck{structures} \add{relations} $\leq$, $\ll$, $\lessgtr$, $\parallel$ sur $\mathcal{F}_{\mathcal{M}}$ induisent des \struck{structures} \add{relations} de même nom sur $\mathcal{F}$.
Les prop~$1_0F$, $2_0F$ (p.~2bis), prop.~2, 3, 5 sont « ipso facto » encore vraies dans un atelier général. Il n'y a que la prop.~4 (recollement des raffinements) qui, dans le cas d'un atelier, demande un axiome \uncertain{additionnel} \emph{\ill{}}, de stabilité, \uncertain{nécessaire} pour l'atelier,

\emph{At Fil} \add{\uncertain{savoir}} Soit $F \in \mathcal{F}$, et soit $(G_\alpha)$ une famille filtrante croissante \struck{pour $\leq$} \add{pour $\leq$} de raffinements de $F$. Supposons que $\forall X \in \widetilde{F}$, la famille des $G_\alpha | X$ soit stationnaire. Alors $\mathrm{Sup}\, G_\alpha$ existe dans $\mathcal{F}$ \struck{\ill{}}.

\uncertain{Comme} on a $\mathcal{F} \subset \mathcal{F}_{\mathcal{M}}$, cela signifie donc \uncertain{que} \ill{} (\uncertain{le Sup dans $\mathcal{F}_{\mathcal{M}}$}) est par $\bigcup G_\alpha$ (qui est le Sup dans $\mathcal{F}_{\mathcal{M}}$) est en fait dans $\mathcal{F}$ : propriété de stabilité de $\mathcal{F}$ en tant que partie \struck{\ill{}} fermée de $\mathcal{F}_{\mathcal{M}}, \leq$.

On peut considérer la structure d'atelier comme une donnée indépendante \ill{} de celle d'un magasin \ill{} \ill{} près.

\emph{Atelier} : $\mathcal{F}$ muni de trois relations $\leq$, $\ll$, $|\circ|$ \ill{}

\emph{At~1} Les Sup majorés existent, et $\mathcal{F} \neq \emptyset$ ($\Leftrightarrow \mathcal{F}$ a un plus petit élément $\emptyset_{\mathcal{F}}$).

\struck{Éléments} Figures élémentaires \uncertain{ou éléments co-irréductibles} :
\[
  X = \mathrm{Sup}\, F_i \Longrightarrow \exists\, i,\ F_i = X .
\]
\struck{\ill{}} $\mathcal{M}$ ens.\ des fig.\ élém., $\widetilde{F} = \{ X \in \mathcal{M} \mid X \leq F \}$.

\emph{At~2} $F \leq G \Longleftrightarrow \widetilde{F} \subset \widetilde{G}$ \struck{\ill{}} … $X \in \mathcal{M}$, $X \leq F \Rightarrow X \leq G$.
« toute strate de $F$ est strate de $G$ ».
Équivalent : Toute figure est le Sup de ses strates.

\marginal{At~2,~3,~4 : p.~55 \uncertain{pour strates}}

\page{20}
\note{il numérote « 12 ».}
\emph{Déf} $F \lessgtr G \Longleftrightarrow \{F, G\}$ majoré ($\Longleftrightarrow F \vee G$ existe)

\emph{At~3} $F \lessgtr G \Longleftrightarrow \forall X \in \widetilde{F},\ Y \in \widetilde{G},\ X \lessgtr Y$

(NB $\Rightarrow$ est déjà \ill{})

\emph{At~4} $F \ll G \Longleftrightarrow \forall X \in \widetilde{F},\ \exists\, Y \in \widetilde{G},\ X \ll Y$

\marginal{NB At~1 -- At~3 ne concernent que $\leq$, At~4 \uncertain{ne fait que préciser} $\ll$ en termes \uncertain{des} strates}

\emph{At~5} a) $F |\circ| G \Longrightarrow F, G \in \mathcal{M}$. \add{i.e.\ la relation $|\circ|$ est en fait une relation dans $\mathcal{M}$}

b) $X \lessgtr Y \Longleftrightarrow \forall X' \in \widetilde{X} \smallsetminus \widetilde{X} \cap \widetilde{Y},\ Y' \in \widetilde{Y} \smallsetminus \widetilde{X} \cap \widetilde{Y}$, on a $X' |\circ| Y'$

\marginal{\uncertain{relation} $|\circ|$ et $\lessgtr$}

$+$ Mag~1, Mag~2, Mag~3, Mag~4

\emph{Scholie} Les deux \uncertain{points de vue} \ill{}

(i) $(\mathcal{M}, \leq, \ll, |\circ|, \mathcal{F} \subset \mathcal{F}_{\mathcal{M}})$ \ill{} \emph{Mag~1 -- Mag~4} $+$ \add{\uncertain{Mag at}}
\note{sous l'accolade de Mag~1 -- Mag~4 : « \ill{} $\leq$, $\ll$, $|\circ|$ ».}

(ii) $(\mathcal{F}, \leq, \ll, |\circ|)$ \ill{} axiomes At~1 -- At~5 $+$ Mag~1 -- Mag~4

sont équivalents, i.e.\ donnent des catégories-groupoïdes équivalentes.

Le point de vue « magasin » est plus pratique, visiblement, pour l'économie de l'écriture des axiomes de base.

Le point de vue « atelier » \uncertain{me semble} plus commode pour la pratique du calcul et des raisonnements géométriques, une fois l'axiomatique acquise.

\note{la page s'arrête ici ; la comparaison des deux points de vue et le programme de la page 13 (lieux, cas ensembliste, subdivisions) se poursuivent au-delà de ce lot.}

\end{document}
